\section{My role}

\subsection*{1a. Roles and responsibilities --- allocated, implicit, and drifted}

On paper I was the autonomy and computer-vision engineer. In practice I became the person who built the parallel universe our project lived in for two months. The explicit allocation --- written down at our kickoff around [PLACEHOLDER: week 13] --- gave me the perception stack: camera pipeline, model training, detection-to-GPS projection, and the mission state machine. Robin owned the flight script and telemetry, Demetro handled presentation and client-facing materials, Edward took the Cube and hardware integration, and [PLACEHOLDER: hardware-lead name] managed the airframe build.

That allocation did not survive contact with reality. When the first flight day in wk18 cancelled on weather, the second in wk19 collapsed on a non-functional airframe, and the third in wk21 on the same, the work I actually did bent hard. I ended up owning the SITL environment, writing the progressive test ladder (\texttt{tests/flight/0a}--\texttt{4}), building \texttt{simple\_simulator.py} (2508 lines), and eventually drafting the team's complaint letter to our course organiser. None of that was in my job description. I took it on because nobody else was going to and because the alternative was two months of dead time. Looking back I can see the drift was not an accident: every time a hardware deadline slipped, the bottleneck shifted upstream to something I could actually build on a laptop, and I followed the bottleneck. That is a useful instinct in a crisis but it is also how I ended up doing four jobs instead of one.

The most honest framing is that I was the \emph{simulation-first} engineer by necessity rather than by design. I set out to do CV. I finished the project having written 2500 lines of simulator, 900 lines of state machine, 752 lines of passive-watch tooling, 127 unit tests, and a 2000-word complaint letter --- and the retrained YOLOv8 model (mAP50 = 0.995) that I had originally considered my ``main'' contribution was in many ways the easiest piece of the lot.

\subsection*{1b. What I am most proud of}

I am most proud of two things, and neither is the one I expected at the start.

The first is the \textbf{simulation stack}. When the wk18 flight cancelled I did not go home --- I sat down and started building the thing that let us develop without a drone. \texttt{simple\_simulator.py} is the piece I keep coming back to in my head. It is a keyboard-driven MVP that flies a virtual drone over the Fenswood map, runs real TFLite inference on synthetic imagery, simulates GPS drift and camera shake, and exercises the full GPS-estimation pipeline from detection pixels to a lat/lon lock. It is ugly in places --- the 2508-line file violates every modularity rule I normally hold myself to --- but it earned its keep. Every bug in the tilt-compensation ray tracing, the inverse-variance weighting, the Kalman filter, and the SSSI-aware landing offset was caught in that simulator before it touched a Pi. When I write in the complaint letter that ``the simulation framework is the only reason this project produced any verifiable engineering output at all,'' I mean it as a factual claim, not a boast. Without it, wk18--wk22 were going to be eight weeks of waiting.

What I am proud of is not that the simulator is clever. It is that I made the call to build it \emph{on the day the flight died}, when the emotionally easier option was to blame the hardware and wait. I now see that the instinct to convert a blocked external dependency into a buildable internal one is the single most useful habit I took away from this module, and I want to keep it.

The second is the \textbf{complaint letter} (\texttt{docs/COMPLAINT\_LETTER.md}). Writing it was uncomfortable. I went through six drafts on the tone --- the early versions were angry in ways that would have been satisfying for me and disastrous for the team. The final version is restrained, factually chronological, names no individuals, and ends with constructive recommendations for future cohorts rather than demands. I sent it on behalf of the team after three cancelled flight days and roughly eight weeks with zero real flight data. The uncomfortable realisation I had while drafting it was that my first instinct --- to absorb the frustration quietly and keep coding --- was not actually professionalism. It was conflict-avoidance dressed up as professionalism, and it was costing the whole team a fair assessment. Writing the letter was the hardest thing I did in this project, harder than any of the code, because it required me to stop hiding in the IDE. I am proud of it because I now think engineering professionalism includes the skill of telling institutions that they have failed you, without burning the bridge.

\subsection*{1c. What I would focus on differently next time}

Less soloing. I will say the uncomfortable thing plainly: I defaulted to solving problems alone under stress, and this project is the first one large enough to reveal that as a pattern rather than a strength. In the wk16--wk18 window, when the hardware was slipping, I wrote three 500-line scripts by myself over weekends instead of asking Robin or Edward to pair on any of them. The code was fine. The loss was that nobody on the team except me understood the state machine well enough to maintain it, and when Robin later needed to integrate his flight script with my detection output in wk20 I had to write him a dedicated integration path (\texttt{passive\_watch\_2.py --robin}, commit \texttt{a5b1ab7}) and a bespoke README (commit \texttt{885cee4}) because he could not read the main codebase fluently. That was a self-inflicted integration tax. Had I paired with Robin on the state machine even two hours a week from wk14 onward, that tax would not have existed.

The second thing I want to do differently is to \textbf{teach} rather than to \textbf{write}. I produced a lot of documentation for teammates --- \texttt{docs/VISION\_PIPELINE\_FOR\_COLLEAGUE.md}, \texttt{COLLEAGUE\_CHECKLIST.md}, the Robin-specific README, the blueprint files under \texttt{docs/}. In retrospect I over-invested in written handoffs and under-invested in live knowledge transfer. A 30-minute whiteboard session with Robin in wk17 would have done more for our integration than any of the markdown I wrote. I now see that my preference for writing over talking is a cost-avoidance move --- writing is a solo activity and I am good at it, whereas a whiteboard session requires me to manage another person's attention in real time, which I find harder. Noticing that is uncomfortable. Fixing it is the single specific CPD commitment I am taking into my next project: for every long doc I write, I will schedule a live walkthrough with the intended reader before I consider the handoff done.
